\documentclass[11pt]{article}
\usepackage[margin=1in]{geometry}
\usepackage{amsmath,amssymb,amsthm}
\usepackage{graphicx}
\usepackage[hidelinks]{hyperref}

\newtheorem{theorem}{Theorem}
\newtheorem{lemma}{Lemma}
\newtheorem{corollary}{Corollary}
\newtheorem{remark}{Remark}

\newcommand{\JF}{\mathrm{JF}}
\newcommand{\dist}{\mathrm{dist}}
\newcommand{\R}{\mathbb{R}}
\newcommand{\1}{\mathbf{1}}

\title{A Partial Answer for Ball Families in the James Function Space Problem}
\author{Run \texttt{fa\_banach\_001}, lane 0}
\date{2026-06-23}

\begin{document}
\maketitle

\begin{center}
\fbox{\parbox{0.9\linewidth}{
\textbf{Status.} Partial result, likely valid.  The packet proves the
nonseparability-of-the-dual part of Problem 1 in arXiv:1210.2379 for the
Euclidean-ball and $\ell_\infty^{d_0}$-ball families.  It does not prove the
remaining assertion that the corresponding spaces do not contain $\ell_1$.
}}
\end{center}

\section{Source Problem}

The source paper is S. A. Argyros, A. Manoussakis, and M. Petrakis,
\emph{Function spaces not containing $\ell_1$}, arXiv:1210.2379
\cite{AMP}.  In Section ``Remarks and Problems,'' Problem 1 asks the following.
Let $\Omega$ be an open bounded subset of $\R^{d_0}$, $d_0>1$, and for a family
$Q(\Omega)$ of convex bodies contained in $\Omega$ define
\[
\|f\|_{\JF(Q(\Omega))}
=\sup\left\{\left(\sum_{i=1}^n\left(\int_{V_i}f\right)^2\right)^{1/2}:
\{V_i\}_{i=1}^n\subset Q(\Omega),\ V_i\cap V_j=\emptyset\right\}.
\]
The problem asks for which families $Q(\Omega)$ the following hold:
\[
\text{(i) } \JF(Q(\Omega))\hbox{ does not contain }\ell_1,\qquad
\text{(ii) } \JF^*(Q(\Omega))\hbox{ is nonseparable.}
\]
The paper notes that this was open, in particular, when $Q(\Omega)$ is the
family of Euclidean balls or the family of $\ell_\infty^{d_0}$ balls contained
in $\Omega$.

\begin{figure}[h]
\centering
\includegraphics[width=\linewidth]{figures/open_problem_crop.png}
\caption{Problem 1 and the ball-family remark in arXiv:1210.2379, PDF page 45.}
\end{figure}

\section{Proof Intuition}

The dual-nonseparability half has a simple mechanism that does not require the
delicate weak-Cauchy analysis used in the paper for excluding $\ell_1$.  Every
set $V$ in the defining family gives an integration functional
$V^*(f)=\int_V f$, and normalized indicators provide exact norm tests for
these functionals.  If two sets $U,V$ differ by positive measure, then the
normalized indicator of $U\setminus V$ (or $V\setminus U$) separates
$U^*$ from $V^*$ by distance at least one.  Thus any uncountable family of
sets that are pairwise distinct modulo null sets yields an uncountable
1-separated subset of the dual.

For Euclidean balls and $\ell_\infty$ balls, such an uncountable family is
obtained by fixing a center in $\Omega$ and varying the radius through an
interval small enough that all the balls remain in $\Omega$.  This proves part
(ii) for the ball families singled out in the source problem.  The argument
does not address part (i), the absence of $\ell_1$.

\section{Statement}

\begin{theorem}\label{thm:uncountable-separated}
Let $\Omega$ be a bounded open subset of $\R^d$ with Lebesgue measure $\mu$,
and let $Q$ be a family of positive-measure measurable subsets of $\Omega$.
Assume that the displayed formula above defines the norm used to form
$\JF(Q)$ from $L^1(\Omega)$, and for $V\in Q$ let
$V^*(f)=\int_V f$ on the dense copy of $L^1(\Omega)$ in $\JF(Q)$.

If $Q$ contains an uncountable subfamily $\mathcal U$ such that any two
distinct $U,V\in\mathcal U$ differ modulo null sets, then $\JF(Q)^*$ is
nonseparable.  More precisely, $\{V^*:V\in\mathcal U\}$ is an uncountable
1-separated subset of $\JF(Q)^*$.
\end{theorem}

\begin{corollary}\label{cor:balls}
Let $\Omega$ be a nonempty bounded open subset of $\R^d$, $d>1$.  If
$Q(\Omega)$ is either the family of Euclidean balls contained in $\Omega$ or
the family of $\ell_\infty^d$ balls contained in $\Omega$, then
$\JF^*(Q(\Omega))$ is nonseparable.
\end{corollary}

\section{Proof}

\begin{lemma}\label{lem:indicator-bound}
Let $A\subset\Omega$ be measurable with $0<\mu(A)<\infty$, and put
$g_A=\1_A/\mu(A)$.  Then $\|g_A\|_{\JF(Q)}\leq 1$.
\end{lemma}

\begin{proof}
For every finite pairwise disjoint family $V_1,\ldots,V_n$ from $Q$,
\[
\left(\sum_{i=1}^n \left|\int_{V_i}g_A\right|^2\right)^{1/2}
=
\left(\sum_{i=1}^n \left(\frac{\mu(V_i\cap A)}{\mu(A)}\right)^2\right)^{1/2}.
\]
All terms are nonnegative, so this is at most
\[
\sum_{i=1}^n \frac{\mu(V_i\cap A)}{\mu(A)}
\leq \frac{\mu(A)}{\mu(A)}=1,
\]
because the sets $V_i\cap A$ are pairwise disjoint subsets of $A$.  Taking the
supremum over all such families gives the claim.
\end{proof}

\begin{proof}[Proof of Theorem \ref{thm:uncountable-separated}]
For each $V\in Q$, the functional $V^*$ is bounded with norm at most one, since
the singleton family $\{V\}$ is permitted in the definition of the norm:
\[
\left|\int_V f\right|\leq \|f\|_{\JF(Q)}.
\]
The bound is sharp.  By Lemma \ref{lem:indicator-bound},
$\|\1_V/\mu(V)\|_{\JF(Q)}\leq 1$, while
\[
V^*(\1_V/\mu(V))=1.
\]
Thus $\|V^*\|=1$.

Now let $U,V\in\mathcal U$ be distinct.  Since they are distinct modulo null
sets, either $\mu(U\setminus V)>0$ or $\mu(V\setminus U)>0$.  In the first
case set $A=U\setminus V$; in the second set $A=V\setminus U$.  Suppose for
definiteness that $\mu(U\setminus V)>0$.  Lemma \ref{lem:indicator-bound}
gives $\|g_A\|_{\JF(Q)}\leq 1$, and
\[
(U^*-V^*)(g_A)
=\int_U g_A-\int_V g_A
=1-0=1.
\]
Hence $\|U^*-V^*\|\geq 1$.  The other case is identical after reversing the
roles of $U$ and $V$.  Therefore $\{V^*:V\in\mathcal U\}$ is uncountable and
1-separated.  A separable metric space cannot contain an uncountable
1-separated subset, so $\JF(Q)^*$ is nonseparable.
\end{proof}

\begin{proof}[Proof of Corollary \ref{cor:balls}]
Choose $x_0\in\Omega$ and $R>0$ such that the closed Euclidean ball
$\overline B_2(x_0,R)$ is contained in $\Omega$.

For Euclidean balls, the family
\[
\{\overline B_2(x_0,r):0<r<R\}
\]
has cardinality continuum, and two balls with distinct radii differ by a set
of positive Lebesgue measure.  Theorem \ref{thm:uncountable-separated} applies.

For $\ell_\infty^d$ balls, choose $0<r<R/\sqrt d$ and use the closed cubes
\[
\overline B_\infty(x_0,r)=x_0+[-r,r]^d.
\]
The family $\{\overline B_\infty(x_0,r):0<r<R/\sqrt d\}$ is again uncountable
and pairwise distinct modulo null sets, and each cube is contained in
$\overline B_2(x_0,R)\subset\Omega$.  Theorem
\ref{thm:uncountable-separated} gives nonseparability of the dual.

The same conclusion holds if the convention in $Q(\Omega)$ uses open balls
instead of closed convex bodies, since the boundary has Lebesgue measure zero
and the mod-null separation argument is unchanged.
\end{proof}

\section{Verification and Limitations}

The proof is a direct estimate from the defining norm and uses no computational
or external structural theorem.  The core inequality is the elementary estimate
\[
\left(\sum_i a_i^2\right)^{1/2}\leq \sum_i a_i
\]
for nonnegative $a_i$, applied to the measures of the disjoint intersections
$V_i\cap A$.

The result is limited to Problem 1(ii), nonseparability of
$\JF^*(Q(\Omega))$.  It does not establish Problem 1(i), namely that the
corresponding space $\JF(Q(\Omega))$ does not contain $\ell_1$.

Cheap run-index searches for arXiv:1210.2379 and the core terms
``Function spaces not containing $\ell_1$,'' ``Euclidean balls,'' and
``$\ell_\infty$ balls'' did not find an existing packet for this target.
Bounded web searches for the same phrases and arXiv id found no obvious later
paper explicitly resolving the ball-family version of Problem 1.  This is
therefore promoted as a partial theorem for human review, not as a claimed
full solution of the source problem.

\section{Human Review Recommendation}

Reviewers should check that the intended definition of $Q(\Omega)$ in the
source paper permits the chosen ball convention and that the completion of
$L^1(\Omega)$ under the displayed norm is the intended ambient space.  The
separation estimate itself is elementary; the main scope issue is that the
argument gives no information about the $\ell_1$-noncontainment half.

\begin{thebibliography}{9}
\bibitem{AMP}
S. A. Argyros, A. Manoussakis, and M. Petrakis,
\emph{Function spaces not containing $\ell_1$},
arXiv:1210.2379, 2012.
\end{thebibliography}

\end{document}
